% Options for packages loaded elsewhere
\PassOptionsToPackage{unicode}{hyperref}
\PassOptionsToPackage{hyphens}{url}
\documentclass[
]{article}
\usepackage{xcolor}
\usepackage{amsmath,amssymb}
\setcounter{secnumdepth}{-\maxdimen} % remove section numbering
\usepackage{iftex}
\ifPDFTeX
  \usepackage[T1]{fontenc}
  \usepackage[utf8]{inputenc}
  \usepackage{textcomp} % provide euro and other symbols
\else % if luatex or xetex
  \usepackage{unicode-math} % this also loads fontspec
  \defaultfontfeatures{Scale=MatchLowercase}
  \defaultfontfeatures[\rmfamily]{Ligatures=TeX,Scale=1}
\fi
\usepackage{lmodern}
\ifPDFTeX\else
  % xetex/luatex font selection
\fi
% Use upquote if available, for straight quotes in verbatim environments
\IfFileExists{upquote.sty}{\usepackage{upquote}}{}
\IfFileExists{microtype.sty}{% use microtype if available
  \usepackage[]{microtype}
  \UseMicrotypeSet[protrusion]{basicmath} % disable protrusion for tt fonts
}{}
\makeatletter
\@ifundefined{KOMAClassName}{% if non-KOMA class
  \IfFileExists{parskip.sty}{%
    \usepackage{parskip}
  }{% else
    \setlength{\parindent}{0pt}
    \setlength{\parskip}{6pt plus 2pt minus 1pt}}
}{% if KOMA class
  \KOMAoptions{parskip=half}}
\makeatother
\setlength{\emergencystretch}{3em} % prevent overfull lines
\providecommand{\tightlist}{%
  \setlength{\itemsep}{0pt}\setlength{\parskip}{0pt}}
\usepackage{bookmark}
\IfFileExists{xurl.sty}{\usepackage{xurl}}{} % add URL line breaks if available
\urlstyle{same}
\hypersetup{
  hidelinks,
  pdfcreator={LaTeX via pandoc}}

\author{}
\date{}

\begin{document}

\section{Comprehensive Guide to the Extended Kalman Filter
(EKF)}\label{comprehensive-guide-to-the-extended-kalman-filter-ekf}

Imagine you are on a commercial airplane flying through a thick storm.
The pilots cannot see outside, and GPS signals can occasionally become
weak or blocked. How does the autopilot system still track the plane's
exact position, speed, and orientation? It uses math to blend its last
known location with noisy data from the plane's internal sensors. This
magic blending trick is called a \textbf{Kalman Filter}.

In the real world, airplanes don't just move in perfectly straight
lines---they turn, climb, descend, and move in complex curves. To track
these curved movements, we use a special upgraded version called the
\textbf{Extended Kalman Filter (EKF)}.

To understand how this works, we will build the math from the ground up,
starting with simple high school algebra, moving through calculus, and
finishing with advanced linear algebra. We will weave all the
foundational pieces into one smooth story.

\begin{center}\rule{0.5\linewidth}{0.5pt}\end{center}

\subsection{Part 1: The Foundations of
Measurement}\label{part-1-the-foundations-of-measurement}

Before we can build a filter, we must understand how to measure things
and how to measure our own uncertainty.

\subsubsection{1. Variables in Multiple Dimensions (Matrices and
Vectors)}\label{variables-in-multiple-dimensions-matrices-and-vectors}

If we only track a car's forward position, we use a single number (a
scalar) like \(x = 5\). But cars have speed, direction, and 2D
coordinates. * \textbf{Algebra Format:} A single equation:
\(y = 3 \cdot x\) * \textbf{Calculus Format:} A velocity vector is the
derivative of a position vector: \(v(t) = [x'(t), y'(t)]\). *
\textbf{Advanced Math (Linear Algebra) Format:} We group multiple
variables into a list called a \textbf{vector}. We use a grid of numbers
called a \textbf{matrix} to manipulate them all at once. Instead of
\(y = 3x\), we write \(Y = A \cdot X\). We multiply the rows of matrix
\(A\) by the column of vector \(X\) to update all our variables
simultaneously.

\subsubsection{2. Measuring Uncertainty (Variance and
Covariance)}\label{measuring-uncertainty-variance-and-covariance}

Sensors aren't perfect. We need a way to measure how much we
\emph{don't} know. * \textbf{Algebra Format:} \textbf{Variance} measures
how spread out a single variable is from its average (\(\mu\)).
\(Var(x) = \text{Average of } (x - \mu)^2\). \textbf{Covariance}
measures how two variables move together. If position and speed increase
together, their covariance is positive.
\(Cov(x, y) = \text{Average of } ((x - \mu_x) \cdot (y - \mu_y))\). *
\textbf{Calculus Format:} For continuous probability curves, we use
integrals to find variance: \(Var(x) = \int (x - \mu)^2 f(x) dx\) *
\textbf{Advanced Math (Linear Algebra) Format:} We track the uncertainty
of an entire vector using a \textbf{Covariance Matrix} (\(\Sigma\) or
\(P\)). It holds the variance of each variable on the diagonal and the
covariances on the off-diagonals.
\(\Sigma = \begin{bmatrix} Var(x) & Cov(x, y) \\ Cov(y, x) & Var(y) \end{bmatrix}\)

\subsubsection{3. The Shape of Probability (Normal
Distribution)}\label{the-shape-of-probability-normal-distribution}

When sensor errors happen, they usually cluster around the true value in
a bell-shaped curve. * \textbf{Algebra Format:} A 1D bell curve defined
by a center (mean \(\mu\)) and a width (standard deviation \(\sigma\)).
* \textbf{Calculus Format:} The exact probability density function:
\(f(x) = \frac{1}{\sigma \sqrt{2\pi}} e^{-\frac{1}{2}(\frac{x-\mu}{\sigma})^2}\)
* \textbf{Advanced Math (Linear Algebra) Format:} A 3D probability cloud
(Multivariate Normal Distribution). We replace division by variance
\(\sigma^2\) with multiplication by the inverse of the covariance matrix
\(\Sigma^{-1}\).
\(f(X) \propto \exp\left( -\frac{1}{2} (X - \mu)^T \Sigma^{-1} (X - \mu) \right)\)

\begin{center}\rule{0.5\linewidth}{0.5pt}\end{center}

\subsection{Part 2: The Standard Linear Kalman
Filter}\label{part-2-the-standard-linear-kalman-filter}

Now we have the tools. Let's derive the standard Kalman Filter, which
assumes everything moves in perfect straight lines. It happens in two
steps: Predict and Correct.

\subsubsection{Step 1: The Prediction}\label{step-1-the-prediction}

We guess where the car is based on its last state.

\begin{itemize}
\item
  \textbf{Algebra Format (1D):} New Position = (Multiplier \(\cdot\) Old
  Position) + Noise \(x_{new} = a \cdot x_{old} + w\) Uncertainty
  expands when we predict: \(p_{new} = a^2 \cdot p_{old} + q\) (where
  \(q\) is process noise).
\item
  \textbf{Calculus Format:} Using derivatives, if position
  \(x(t) = v \cdot t\), the change over time is \(\frac{dx}{dt} = v\).
  We use this constant rate to step forward in time.
\item
  \textbf{Advanced Math (Linear Algebra) Format:} State Prediction:
  \(X_{new} = A \cdot X_{old} + W\) Uncertainty Prediction: You can't
  just ``square'' a matrix \(A\) like we squared \(a\) in algebra. You
  multiply it by its transpose!
  \(P_{new} = A \cdot P_{old} \cdot A^T + Q\)
\end{itemize}

\subsubsection{Step 2: The Correction}\label{step-2-the-correction}

We read the GPS sensor. It gives us a measurement \(Z\). The sensor
reading relates to our true state by a multiplier \(H\). How much do we
trust the sensor versus our prediction? We calculate a ratio called the
\textbf{Kalman Gain (\(K\))}.

\begin{itemize}
\item
  \textbf{Algebra Format (1D):}
  \(K = \frac{\text{Prediction Uncertainty}}{\text{Prediction Uncertainty} + \text{Sensor Noise (R)}}\)
  \(K = \frac{p \cdot h}{h \cdot p \cdot h + R}\) Final State =
  Prediction +
  \(K \cdot (\text{Actual Sensor} - \text{Expected Sensor})\)
  \(x_{final} = x_{new} + K \cdot (z - h \cdot x_{new})\)
\item
  \textbf{Advanced Math (Linear Algebra) Format:} Kalman Gain (matrix
  division is multiplying by the inverse):
  \(K = P_{new} \cdot H^T \cdot (H \cdot P_{new} \cdot H^T + R)^{-1}\)

  Final State Update:
  \(X_{final} = X_{new} + K \cdot (Z - H \cdot X_{new})\)

  Final Uncertainty Update (our uncertainty shrinks because we added new
  info!): \(P_{final} = (I - K \cdot H) \cdot P_{new}\) (where \(I\) is
  the identity matrix, acting like the number 1).
\end{itemize}

\begin{center}\rule{0.5\linewidth}{0.5pt}\end{center}

\subsection{Part 3: Handling Curves (The EKF
Upgrade)}\label{part-3-handling-curves-the-ekf-upgrade}

The standard Kalman filter uses fixed matrices \(A\) and \(H\). This
strictly represents straight lines. But if our car turns, the math
function \(f(x)\) becomes curvy (non-linear). To use our Kalman
equations, we must straighten out the curves!

\subsubsection{1. Slopes and Partial
Derivatives}\label{slopes-and-partial-derivatives}

To straighten a curve, we need to know its slope. * \textbf{Algebra
Format:} Slope is rise over run. \(m = \frac{y_2 - y_1}{x_2 - x_1}\) *
\textbf{Calculus Format:} The derivative gives the exact slope at a
point: \(f'(x)\). A \textbf{partial derivative}
\(\frac{\partial f}{\partial x}\) gives the slope in one specific
direction.

\subsubsection{2. Straightening the Curve (Taylor
Series)}\label{straightening-the-curve-taylor-series}

\begin{itemize}
\tightlist
\item
  \textbf{Algebra Format:} We estimate a nearby point on a curve by
  drawing a straight line:
  \(y_{est} = y_{start} + m \cdot (x - x_{start})\).
\item
  \textbf{Calculus Format:} The Taylor Series proves this.
  \(f(x) = f(a) + f'(a)(x-a) + \frac{f''(a)}{2}(x-a)^2 + \dots\) If we
  zoom in extremely close, \((x-a)\) is tiny. Squaring a tiny number
  makes it effectively zero. We can delete the rest of the equation! We
  are left with a perfect linear approximation:
  \(f(x) \approx f(a) + f'(a)(x-a)\)
\end{itemize}

\subsubsection{3. The Multi-Variable Slope (The Jacobian
Matrix)}\label{the-multi-variable-slope-the-jacobian-matrix}

Because our car has multiple variables, our derivative \(f'(a)\) isn't a
single number. We take all the partial derivatives of all our variables
and pack them into a grid called the \textbf{Jacobian Matrix (\(J\))}. *
\textbf{Linear Algebra Format:} The linear approximation of a curvy
vector function becomes: \(F(X) \approx F(A) + J \cdot (X - A)\)

\begin{center}\rule{0.5\linewidth}{0.5pt}\end{center}

\subsection{Part 4: Deriving the Extended Kalman
Filter}\label{part-4-deriving-the-extended-kalman-filter}

We now take our standard Kalman filter equations and apply our new
Taylor Series Jacobian trick to them line by line!

\subsubsection{The Prediction Phase
(EKF)}\label{the-prediction-phase-ekf}

\begin{itemize}
\item
  \textbf{Algebra Format:} Instead of a straight line \(y = ax\), our
  state evolves by a curvy function \(f(x)\). State:
  \(x_{new} = f(x_{old})\) Uncertainty: We use the slope of the curve
  \(f'(x)\). \(p_{new} = (f'(x))^2 \cdot p_{old} + q\)
\item
  \textbf{Calculus Format:} By the 1st-order Taylor Series, we
  linearized \(f(x)\) using its derivative \(f'(x)\) evaluated at our
  last best guess. This derivative acts as our new, temporary
  straight-line multiplier.
\item
  \textbf{Advanced Math (Linear Algebra) Format:} State: We just pass
  our vector through the non-linear math function:
  \(X_{new} = f(X_{old})\) Uncertainty: We replace the fixed matrix
  \(A\) with the \textbf{Jacobian Matrix (\(F_k\))} calculated at the
  current moment. \(P_{new} = F_k \cdot P_{old} \cdot F_k^T + Q\)
\end{itemize}

\subsubsection{The Correction Phase
(EKF)}\label{the-correction-phase-ekf}

\begin{itemize}
\item
  \textbf{Algebra Format:} Our sensor reading is also a curvy function
  \(h(x)\). We use its slope \(h'(x)\) to calculate the Kalman gain.
  \(K = \frac{p \cdot h'(x)}{h'(x) \cdot p \cdot h'(x) + R}\)
\item
  \textbf{Calculus Format:} Again, using the Taylor Series, we
  approximate the curvy sensor function \(h(x)\) using its first
  derivative \(h'(x)\). We evaluate this derivative at our newly
  predicted state.
\item
  \textbf{Advanced Math (Linear Algebra) Format:} We calculate the
  Jacobian Matrix of the sensor function, calling it \textbf{\(H_k\)}.

  \textbf{Step-by-Step EKF Update:}

  \begin{enumerate}
  \def\labelenumi{\arabic{enumi}.}
  \tightlist
  \item
    Calculate Kalman Gain using the Jacobian \(H_k\):
    \(K = P_{new} \cdot H_k^T \cdot (H_k \cdot P_{new} \cdot H_k^T + R)^{-1}\)
  \item
    Update State (Note we use the true curvy function \(h(X)\) to find
    the expected sensor reading, not the Jacobian!):
    \(X_{final} = X_{new} + K \cdot (Z - h(X_{new}))\)
  \item
    Update Uncertainty: \(P_{final} = (I - K \cdot H_k) \cdot P_{new}\)
  \end{enumerate}
\end{itemize}

\subsection{Summary of the Journey}\label{summary-of-the-journey}

We started with basic slopes and single variables. We realized that to
track complex things like cars, we needed Matrices and Covariance. We
derived the standard Kalman Filter for perfectly straight movements.
Then, by using Calculus (Derivatives and Taylor Series) to zoom in and
straighten out curves, we upgraded our math to Linear Algebra (Jacobian
Matrices). This birthed the Extended Kalman Filter, allowing us to
perfectly blend noisy predictions and noisy sensors in a completely
curved, non-linear world!

\end{document}
